% Chapter 4, Section 1 _Linear Algebra_ Jim Hefferon
%  http://joshua.smcvt.edu/linearalgebra
%  2001-Jun-12
\section{Nilpotensi}
Bab ini menunjukkan bahwa setiap matriks persegi serupa dengan suatu matriks
yang merupakan jumlah dari dua jenis matriks sederhana.
Bagian sebelumnya berfokus pada jenis sederhana yang pertama, yaitu matriks
diagonal.
Sekarang kita meninjau jenis yang lain.









\subsectionoptional{Komposisi Diri}\index{komposisi diri!pemetaan}\index{pemetaan!komposisi diri}\index{transformasi!dikomposisikan dengan dirinya sendiri}\index{komposisi!diri}
% \noindent\textit{This subsection is optional, although it is necessary
% for later material in this section and in the next one.}

Karena transformasi linear $\map{t}{V}{V}$ memiliki domain dan kodomain yang
sama, kita dapat mengomposisikan $t$ dengan dirinya sendiri:
\( t^2=\composed{t}{t} \), 
  dan \( t^3=\composed{t}{\composed{t}{t}} \), 
dan seterusnya.\appendrefs{iterasi fungsi}\spacefactor=1000 %
\begin{center}
  \includegraphics{jc/mp/ch5.1}
\end{center}
Notasi pangkat superskrip $t^j$ bagi iterasi transformasi sesuai dengan notasi
yang telah kita gunakan bagi representasi matriks perseginya, sebab jika
$\rep{t}{B,B}=T$, maka \( \rep{t^j}{B,B}=T^j \).


\begin{example} \label{ex:DerivIter}
Bagi pemetaan turunan \( \map{d/dx}{\polyspace_3}{\polyspace_3} \) yang
diberikan oleh
\begin{equation*}
  a+bx+cx^2+dx^3\xmapsunder{d/dx} b+2cx+3dx^2
\end{equation*}
pangkat keduanya adalah turunan kedua,
\begin{equation*}
  a+bx+cx^2+dx^3\xmapsunder{d^2/dx^2} 2c+6dx 
\end{equation*}
pangkat ketiganya adalah turunan ketiga,
\begin{equation*}
  a+bx+cx^2+dx^3\xmapsunder{d^3/dx^3} 6d 
\end{equation*}
dan setiap pangkat yang lebih tinggi adalah pemetaan nol.
\end{example}

\begin{example}
Transformasi pada ruang $\matspace_{\nbyn{2}}$ yang terdiri atas matriks
$\nbyn{2}$ ini
\begin{equation*}
  \begin{mat}
    a  &b  \\
    c  &d
  \end{mat}
  \mapsunder{t}
  \begin{mat}
    b  &a  \\
    d  &0
  \end{mat}
\end{equation*}
memiliki pangkat kedua berikut
\begin{equation*}
  \begin{mat}
    a  &b  \\
    c  &d
  \end{mat}
  \mapsunder{t^2}
  \begin{mat}
    a  &b  \\
    0  &0
  \end{mat}
\end{equation*}
dan pangkat ketiga berikut.
\begin{equation*}
  \begin{mat}
    a  &b  \\
    c  &d
  \end{mat}
  \mapsunder{t^3}
  \begin{mat}
    b  &a  \\
    0  &0
  \end{mat}
\end{equation*}
Setelah itu, $t^4=t^2$ dan $t^5=t^3$, dan seterusnya.
\end{example}

\begin{example}
Tinjau transformasi geser $\map{t}{\C^3}{\C^3}$.
\begin{equation*}
  \colvec{x \\ y \\ z} \mapsunder{t} \colvec{0 \\ x \\ y}
\end{equation*}
Kita memperoleh
\begin{equation*}
  \colvec{x \\ y \\ z} \mapsunder{t} \colvec{0 \\ x \\ y}
                       \mapsunder{t} \colvec{0 \\ 0 \\ x}
                       \mapsunder{t} \colvec{0 \\ 0 \\ 0}
\end{equation*}
sehingga ruang-ruang hasilnya turun menuju subruang trivial.
\begin{equation*}
  \rangespace{t}=\set{\colvec{0 \\ a \\ b}\suchthat a,b\in\C}
  \qquad
  \rangespace{t^2}=\set{\colvec{0 \\ 0 \\ c}\suchthat c\in\C}
  \qquad
  \rangespace{t^3}=\set{\colvec{0 \\ 0 \\ 0}}  
\end{equation*}
\end{example}

Contoh-contoh ini mengisyaratkan bahwa setelah sejumlah iterasi, pemetaan
tersebut menjadi stabil.

\begin{lemma}  \label{le:RangeAndNullChains}
Bagi sebarang transformasi \( \map{t}{V}{V} \), ruang hasil dari
pangkat-pangkatnya membentuk rantai menurun
\begin{equation*}
  V\supseteq \rangespace{t}\supseteq\rangespace{t^2}\supseteq\cdots
\end{equation*}
dan ruang nolnya membentuk rantai menaik.
\begin{equation*}
  \set{\vec{0}}\subseteq\nullspace{t}\subseteq\nullspace{t^2}\subseteq\cdots
\end{equation*}
Lebih lanjut, terdapat \( k>0 \) sedemikian sehingga bagi pangkat-pangkat yang
lebih kecil daripada $k$, himpunan bagiannya bersifat sejati:~jika
$j<k$, maka $\rangespace{t^j}\supset\rangespace{t^{j+1}}$
dan 
$\nullspace{t^j}\subset\nullspace{t^{j+1}}$,
sedangkan
jika $j\geq k$, maka $\rangespace{t^j}=\rangespace{t^{j+1}}$
dan
$\nullspace{t^j}=\nullspace{t^{j+1}}$.  
\end{lemma}

\noindent (Kasus $k=1$ dapat terjadi, misalnya jika $t$ adalah pemetaan
identitas, sehingga tidak ada himpunan bagian sejati dalam kedua rantai.)

\begin{proof}
Pertama, ingat bahwa bagi sebarang pemetaan, dimensi ruang hasil ditambah dimensi
ruang nol sama dengan dimensi domainnya.
Jadi, jika dimensi ruang hasil berkurang, dimensi ruang nol harus bertambah.
Di sini kita akan membuktikan bagian mengenai ruang hasil dan menyerahkan sisanya
kepada
\nearbyexercise{exer:RangeAndNullChains}.

Kita mulai dengan menunjukkan bahwa ruang-ruang hasil tersebut membentuk suatu
rantai.
Jika $\vec{w}\in\rangespace{t^{j+1}}$, sehingga
$\vec{w}=t^{j+1}(\vec{v})$ bagi suatu $\vec{v}$, maka
$\vec{w}=t^{j}(\,t(\vec{v})\,)$.
Jadi, $\vec{w}\in\rangespace{t^{j}}$.

Berikutnya kita membuktikan pernyataan ``lebih lanjut'':~pada awalnya
inklusi-inklusi dalam rantai bersifat sejati, lalu mulai suatu pangkat $k$
ruang-ruang hasil tersebut sama.
Pertama kita tunjukkan bahwa jika sepasang ruang hasil yang bersebelahan dalam
rantai sama, \( \rangespace{t^{k}}=\rangespace{t^{k+1}} \), maka semua ruang
hasil sesudahnya juga sama:
\( \rangespace{t^{k+1}}=\rangespace{t^{k+2}} \), dan seterusnya.
Hal ini berlaku karena
\( \map{t}{\rangespace{t^{k+1}}}{\rangespace{t^{k+2}}} \)
adalah pemetaan yang sama, dengan domain yang sama, seperti
\( \map{t}{\rangespace{t^{k}}}{\rangespace{t^{k+1}}} \), dan
karena itu mempunyai hasil yang sama,
\( \rangespace{t^{k+1}}=\rangespace{t^{k+2}} \)
(pernyataan ini berlaku bagi semua pangkat yang lebih tinggi melalui induksi).
Jadi, jika rantai ruang hasil berhenti menurun secara ketat, mulai titik itu
rantai tersebut stabil.

Terakhir, kita tunjukkan bahwa rantai tersebut pada akhirnya harus berhenti
menurun.
Setiap ruang hasil merupakan subruang dari ruang sebelumnya.
Agar menjadi subruang sejati, dimensinya harus benar-benar lebih rendah
(lihat \nearbyexercise{exer:PropSubspStrictLowerDimen}).
Ruang-ruang ini berdimensi hingga, sehingga rantai hanya dapat turun sebanyak
berhingga langkah.
Artinya, pangkat $k$ paling besar sama dengan dimensi $V$.
\end{proof}

\begin{example}
Pemetaan turunan $a+bx+cx^2+dx^3\xmapsunder{d/dx} b+2cx+3dx^2$ pada
$\polyspace_3$ memiliki rantai ruang hasil berikut.
\begin{equation*}
  \rangespace{t^0}=\polyspace_3
    \:\supset\:\rangespace{t^1}=\polyspace_2
    \:\supset\:\rangespace{t^2}=\polyspace_1
    \:\supset\:\rangespace{t^3}=\polyspace_0
    \:\supset\:\rangespace{t^4}=\set{\zero}
    % =\rangespace{t^5}=\set{\zero}=
    % \,=\,\cdots
\end{equation*}
Semua unsur rantai sesudahnya adalah ruang trivial.
Pemetaan tersebut memiliki rantai ruang nol berikut.
\begin{equation*}
  \nullspace{t^0}=\set{\zero}
  \:\subset\: \nullspace{t^1}=\polyspace_0
  \:\subset\:\nullspace{t^2}=\polyspace_1
  \:\subset\:\nullspace{t^3}=\polyspace_2
  \:\subset\:\nullspace{t^4}=\polyspace_3
\end{equation*}
Unsur-unsur sesudahnya adalah seluruh ruang.
\end{example}

\begin{example} \label{exam:PolyRankFalls}
Misalkan \( \map{t}{\polyspace_2}{\polyspace_2} \) adalah pemetaan
\( d_0+d_1x+d_2x^2 \mapsto 2d_0+d_2x. \)
Seperti dinyatakan lema, melalui iterasi ruang hasilnya menyusut
\begin{equation*}
  \rangespace{t^0}=\polyspace_2
    \quad
  \rangespace{t}=\set{a_0+a_1x\suchthat a_0,a_1\in\C}
    \quad
  \rangespace{t^2}=\set{a_0\suchthat a_0\in\C}
\end{equation*}
kemudian stabil, sehingga $\rangespace{t^2}=\rangespace{t^3}=\cdots$.
Ruang nolnya membesar
\begin{equation*}
  \nullspace{t^0}=\set{0}
    \quad
  \nullspace{t}=\set{b_1x\suchthat b_1\in\C}
    \quad
  \nullspace{t^2}=\set{b_1x+b_2x^2\suchthat b_1,b_2\in\C}
\end{equation*}
kemudian stabil, $\nullspace{t^2}=\nullspace{t^3}=\cdots$.
\end{example}

\begin{example}
Transformasi \( \map{\pi}{\C^3}{\C^3} \) yang memproyeksikan ke dua koordinat
pertama,
\begin{equation*}
   \colvec{c_1 \\ c_2 \\ c_3}
     \mapsunder{\pi}
   \colvec{c_1 \\ c_2 \\ 0}
\end{equation*}
memiliki \( \C^3\supset\rangespace{\pi}=\rangespace{\pi^2}=\cdots \) dan
\( \set{\zero}\subset\nullspace{\pi}=\nullspace{\pi^2}=\cdots\, \), dengan
ruang hasil dan ruang nol berikut.
\begin{equation*}
  \rangespace{\pi}=\set{\colvec{a \\ b \\ 0}\suchthat a,b\in\C}
  \qquad
  \nullspace{\pi}=\set{\colvec{0 \\ 0 \\ c}\suchthat c\in\C}
\end{equation*}
\end{example}

\begin{definition}
Misalkan \( t \) adalah transformasi pada ruang berdimensi \( n \).
\definend{Ruang hasil tergeneralisasi}\index{ruang hasil tergeneralisasi}%
\index{ruang hasil!tergeneralisasi}
(atau \definend{penutupan ruang hasil\/}\index{penutupan!ruang hasil}%
\index{ruang hasil!penutupan})
adalah $\genrangespace{t}=\rangespace{t^n}$.
\definend{Ruang nol tergeneralisasi}\index{ruang nol tergeneralisasi}%
\index{ruang nol!tergeneralisasi}
(atau \definend{penutupan ruang nol\/}\index{penutupan!ruang nol}%
\index{ruang nol!penutupan})
adalah $\gennullspace{t}=\nullspace{t^n}$.
\end{definition}

Grafik berikut menggambarkannya.
Sumbu horizontal menunjukkan pangkat~$j$ suatu transformasi.
Sumbu vertikal menunjukkan dimensi ruang hasil $t^j$ sebagai jarak di atas nol,
dan dengan demikian juga menunjukkan dimensi ruang nol karena jumlah keduanya
sama dengan dimensi $n$ domain.
\begin{center}
  % \includegraphics{jc/mp/ch5.2}
  \includegraphics{jc/mp/ch5.8}
\end{center}
Melalui iterasi, peringkat turun dan nulitas naik hingga terdapat suatu $k$ yang
membuat pemetaan mencapai keadaan stabil,
$\rangespace{t^k}=\rangespace{t^{k+1}}=\genrangespace{t}$
dan $\nullspace{t^k}=\nullspace{t^{k+1}}=\gennullspace{t}$.
Hal ini harus terjadi paling lambat pada iterasi ke-$n$.
% The steady state's distance above zero is the dimension of the 
% generalized range space
% and its distance below $n$ is the dimension of
% the generalized null space.

\begin{exercises}
  \recommended \item 
    Berikan rantai ruang hasil dan ruang nol bagi transformasi nol dan
    transformasi identitas.
    \begin{answer}
      Bagi transformasi nol, apa pun ruangnya, rantai ruang hasil adalah
      \( V\supset\set{\vec{0}}=\set{\vec{0}}=\cdots\, \), sedangkan rantai
      ruang nol adalah \( \set{\vec{0}}\subset V=V=\cdots\, \).
      Bagi transformasi identitas, kedua rantai tersebut adalah
      \( V=V=V=\cdots \) dan
      \( \set{\vec{0}}=\set{\vec{0}}=\cdots\, \).
    \end{answer}
  \recommended\item 
     Bagi setiap pemetaan, berikan rantai ruang hasil, rantai ruang nol, ruang
     hasil tergeneralisasi, dan ruang nol tergeneralisasi.
     \begin{exparts}
       \partsitem $\map{t_0}{\polyspace_2}{\polyspace_2}$, 
         $a+bx+cx^2\mapsto b+cx^2$ 
       \partsitem $\map{t_1}{\Re^2}{\Re^2}$,
         \begin{equation*}
           \colvec{a \\ b}\mapsto\colvec{0 \\ a}
         \end{equation*}
       \partsitem $\map{t_2}{\polyspace_2}{\polyspace_2}$, 
         $a+bx+cx^2\mapsto b+cx+ax^2$
       \partsitem $\map{t_3}{\Re^3}{\Re^3}$,
         \begin{equation*}
           \colvec{a \\ b \\ c}\mapsto\colvec{a \\ a \\ b}
         \end{equation*}
     \end{exparts}
     \begin{answer}
       \begin{exparts}
         \partsitem Mengiterasikan $t_0$ dua kali,
           $a+bx+cx^2\mapsto b+cx^2\mapsto cx^2$
           memberikan
           \begin{equation*}
             a+bx+cx^2\mapsunder{t_0^2}cx^2
           \end{equation*}
           dan setiap iterasi yang lebih tinggi merupakan pemetaan yang sama.
           Jadi, $\rangespace{t_0}$ adalah ruang polinomial kuadrat tanpa suku
           linear $\set{p+rx^2\suchthat p,r\in \C}$, sedangkan
           $\rangespace{t_0^2}$ adalah ruang polinomial kuadrat murni
           $\set{rx^2\suchthat r\in \C}$; di sinilah rantai menjadi stabil,
           $\genrangespace{t_0}=\set{rx^2\suchthat r\in \C}$.
 
           Bagi ruang nol, $\nullspace{t_0}$ adalah ruang polinomial konstan
           $\set{p\suchthat p\in \C}$, sedangkan $\nullspace{t_0^2}$ adalah
           ruang polinomial linear $\set{p+qx\suchthat p,q\in \C}$.
           Inilah akhirnya: $\gennullspace{t_0}=\nullspace{t_0^2}$.
         \partsitem Pangkat kedua
           \begin{equation*}
             \colvec{a \\ b}
              \mapsunder{t_1}\colvec{0 \\ a}
              \mapsunder{t_1}\colvec{0 \\ 0}
           \end{equation*}
           adalah pemetaan nol.
           Akibatnya, rantai ruang hasil
           \begin{equation*}
             \Re^2
               \supset\set{\colvec{0 \\ p}\suchthat p\in\Re}
               \supset\set{\zero}
               =\cdots
           \end{equation*}
           dan rantai ruang nol
           \begin{equation*}
             \set{\zero}
               \subset\set{\colvec{0 \\ q}\suchthat q\in\Re}
               \subset\Re^2
               =\cdots
           \end{equation*}
           masing-masing memiliki panjang dua.
           Ruang hasil tergeneralisasinya adalah subruang trivial, sedangkan
           ruang nol tergeneralisasinya adalah seluruh ruang.
         \partsitem Iterasi pemetaan ini berputar dalam siklus
           \begin{equation*}
             a+bx+cx^2
               \mapsunder{t_2} b+cx+ax^2            
               \mapsunder{t_2} c+ax+bx^2            
               \mapsunder{t_2} a+bx+cx^2
               \;\cdots            
           \end{equation*}
           dan rantai ruang hasil serta ruang nolnya bersifat trivial.
           \begin{equation*}
             \polyspace_2=\polyspace_2=\cdots
              \qquad
              \set{\zero}=\set{\zero}=\cdots  
           \end{equation*}
           Jadi, jelas bahwa ruang-ruang tergeneralisasinya adalah
           $\genrangespace{t_2}=\polyspace_2$ dan
           $\gennullspace{t_2}=\set{\zero}$.
         \partsitem Kita memperoleh
           \begin{equation*}
             \colvec{a \\ b \\ c}
                \mapsto\colvec{a \\ a \\ b}
                \mapsto\colvec{a \\ a \\ a}
                \mapsto\colvec{a \\ a \\ a}
                \mapsto\cdots
           \end{equation*}
           sehingga rantai ruang hasil
           \begin{equation*}
             \Re^3
               \supset\set{\colvec{p \\ p \\ r}\suchthat p,r\in\Re}
               \supset\set{\colvec{p \\ p \\ p}\suchthat p\in\Re}
               =\cdots
           \end{equation*}
           dan rantai ruang nol
           \begin{equation*}
             \set{\zero}
                 \subset\set{\colvec{0 \\ 0 \\ r}\suchthat r\in \Re}
                 \subset\set{\colvec{0 \\ q \\ r}\suchthat q,r\in \Re}
                =\cdots
           \end{equation*}
           masing-masing memiliki panjang dua.
           Ruang-ruang tergeneralisasinya adalah ruang terakhir yang ditampilkan
           pada masing-masing rantai di atas.
      \end{exparts}
     \end{answer}
  \item 
    Buktikan bahwa komposisi fungsi bersifat asosiatif,
    \( \composed{(\composed{t}{t})}{t}=\composed{t}{(\composed{t}{t})} \)
    sehingga kita dapat menulis $t^3$ tanpa menentukan pengelompokan.
    \begin{answer}
      Keduanya memetakan \( x\mapsto t(t(t(x))) \).
    \end{answer}
  \item \label{exer:PropSubspStrictLowerDimen}
    Periksa bahwa dimensi suatu subruang harus kurang dari atau sama dengan
    dimensi ruang yang memuatnya.
    Periksa bahwa jika subruang itu sejati (subruang tersebut tidak sama dengan
    ruang yang memuatnya), maka dimensinya benar-benar lebih kecil.
    \textit{(Hal ini digunakan dalam bukti
             \nearbylemma{le:RangeAndNullChains}.)}
    \begin{answer}
      Ingat bahwa jika $W$ adalah subruang dari $V$, kita dapat memperluas
      sebarang basis $B_W$ bagi $W$ menjadi basis $B_V$ bagi $V$.
      Pernyataan pertama langsung mengikuti fakta ini.
      Pernyataan kedua juga tidak sulit:~$W$ adalah rentang $B_W$, dan jika $W$
      merupakan subruang sejati, maka $V$ bukan rentang $B_W$. Karena itu,
      $B_V$ harus memiliki sekurang-kurangnya satu vektor lebih banyak daripada
      $B_W$.
    \end{answer}
  \recommended\item 
    Buktikan bahwa ruang hasil tergeneralisasi $\genrangespace{t}$ adalah seluruh
    ruang dan ruang nol tergeneralisasi $\gennullspace{t}$ bersifat trivial jika
    transformasi $t$ nonsingular.
    Apakah pernyataan ``hanya jika'' juga berlaku?
    \begin{answer}
      Pernyataan ``jika'' dan ``hanya jika'' keduanya berlaku.
      Suatu pemetaan linear nonsingular jika dan hanya jika mempertahankan
      dimensi, yaitu jika dimensi hasilnya sama dengan dimensi domainnya.
      Bagi transformasi $\map{t}{V}{V}$, ini berarti pemetaan nonsingular jika
      dan hanya jika surjektif:~$\rangespace{t}=V$ (dan karena itu
      $\rangespace{t^2}=V$, dan seterusnya).
    \end{answer}
  \item \label{exer:RangeAndNullChains} 
    Buktikan bagian mengenai ruang nol dari
    \nearbylemma{le:RangeAndNullChains}.
    \begin{answer}
      Ruang-ruang nol membentuk rantai karena jika
      $\vec{v}\in\nullspace{t^j}$, maka $t^j(\vec{v})=\zero$ dan
      $t^{j+1}(\vec{v})=t(\,t^j(\vec{v})\,)=t(\zero)=\zero$, sehingga
      $\vec{v}\in\nullspace{t^{j+1}}$.

      Selanjutnya, sifat ``lebih lanjut'' bagi ruang nol mengikuti fakta bahwa
      sifat itu berlaku bagi ruang hasil, bersama dengan latihan sebelumnya.
      Karena dimensi $\rangespace{t^j}$ ditambah dimensi
      $\nullspace{t^j}$ sama dengan dimensi~$n$ ruang awal~$V$, ketika dimensi
      ruang hasil berhenti menurun, dimensi ruang nol juga berhenti meningkat.
      Latihan sebelumnya menunjukkan bahwa mulai titik~$k$ ini, inklusi dalam
      rantai tidak lagi sejati\Dash ruang-ruang nol tersebut sama.
    \end{answer}
  \recommended\item
    Berikan contoh transformasi pada ruang berdimensi tiga yang hasilnya
    berdimensi dua.
    Apakah ruang nolnya?
    Iterasikan contoh Anda hingga ruang hasil dan ruang nolnya stabil.
    \begin{answer}
      (Banyak contoh yang benar; berikut salah satunya.)
      Salah satu contoh adalah operator geser pada tripel bilangan real,
      \( (x,y,z)\mapsto (0,x,y) \).
      Ruang nolnya terdiri atas semua tripel yang diawali dua nol.
      Pemetaan tersebut stabil setelah tiga iterasi.
     \end{answer}
  \item 
      Tunjukkan bahwa ruang hasil dan ruang nol suatu transformasi linear tidak
      harus saling lepas.
      Apakah keduanya pernah saling lepas?
      \begin{answer}
        Operator diferensiasi \( \map{d/dx}{\polyspace_1}{\polyspace_1} \)
        memiliki ruang hasil yang sama dengan ruang nolnya.
        Sebagai contoh ketika keduanya saling lepas\Dash kecuali pada vektor
        nol\Dash tinjau pemetaan identitas atau sebarang pemetaan nonsingular.
      \end{answer}
\end{exercises}

















\subsectionoptional{Untaian}
\index{nilpoten|(}
\noindent\textit{Bagian ini memerlukan materi dari subbagian opsional
    Menggabungkan Subruang.}

Subbagian sebelumnya menunjukkan bahwa ketika \( j \) bertambah, dimensi
$\rangespace{t^j}$ turun sedangkan dimensi $\nullspace{t^j}$ naik, sedemikian
rupa sehingga peringkat dan nulitas tersebut bersama-sama membagi dimensi $V$.
Dapatkah kita mengatakan lebih banyak; apakah keduanya membagi suatu basis\Dash
apakah
\( V=\rangespace{t^j}\directsum\nullspace{t^j} \)?

Jawabannya ya bagi pangkat terkecil $j=0$, sebab
\( V=\rangespace{t^0}\directsum\nullspace{t^0}=V\directsum\set{\zero} \).
Jawabannya juga ya pada ujung yang lain.

\begin{lemma} \label{lem:RestONeToOne}
Bagi sebarang transformasi linear \( \map{t}{V}{V} \), fungsi
\( \map{t}{\genrangespace{t}}{\genrangespace{t}} \) bersifat satu-satu.
\end{lemma}

\begin{proof}
Misalkan dimensi $V$ adalah $n$.
Karena \( \rangespace{t^n}=\rangespace{t^{n+1}} \), pemetaan
\( \map{t}{\genrangespace{t}}{\genrangespace{t}} \) merupakan homomorfisme
yang mempertahankan dimensi.
Karena itu, menurut Teorema~Tiga.II.\ref{th:OOHomoEquivalence}, pemetaan tersebut
bersifat satu-satu.
\end{proof}

\begin{corollary} \label{GenRngNullDirSumToSp}
Jika \( \map{t}{V}{V} \) adalah transformasi linear, ruang tersebut merupakan
jumlah langsung
\( V=\genrangespace{t}\directsum\gennullspace{t} \).
Artinya, berlaku (1)~\( \dim(V)=\dim(\genrangespace{t})+\dim(\gennullspace{t}) \)
dan (2)~\( \genrangespace{t}\intersection\gennullspace{t}=\set{\zero} \).
\end{corollary}

\begin{proof}
Misalkan dimensi $V$ adalah $n$.
Kita akan membuktikan kalimat kedua, yang ekuivalen dengan kalimat pertama.
Klausul~(1) benar karena setiap transformasi memenuhi bahwa rank ditambah
nulitasnya sama dengan dimensi ruang, dan secara khusus hal ini berlaku bagi
transformasi $t^n$.

Bagi klausul~(2), andaikan
\( \vec{v}\in\genrangespace{t}\intersection\gennullspace{t} \)
%               =\rangespace{t^n}\intersection\nullspace{t^n} \)
dan kita akan membuktikan bahwa $\vec{v}=\zero$.
Karena \( \vec{v} \) berada dalam ruang nol tergeneralisasi,
\( t^n(\vec{v})=\zero \).
Di sisi lain, menurut lema,
\( \map{t}{\genrangespace{t}}{\genrangespace{t}} \)
bersifat satu-satu. Karena komposisi pemetaan-pemetaan satu-satu juga bersifat
satu-satu, \( \map{t^n}{\genrangespace{t}}{\genrangespace{t}} \) bersifat
satu-satu.
Hanya \( \zero \) yang dipetakan ke \( \zero \) oleh pemetaan linear
satu-satu, sehingga fakta bahwa \( t^n(\vec{v})=\zero \) menyiratkan
\( \vec{v}=\zero \).
\end{proof}

\begin{remark}
Secara teknis terdapat perbedaan antara pemetaan $\map{t}{V}{V}$ dan pemetaan
pada subruang \( \map{t}{\genrangespace{t}}{\genrangespace{t}} \) jika ruang
hasil tergeneralisasi tidak sama dengan $V$, sebab domainnya berbeda.
Namun, perbedaannya kecil karena pemetaan kedua merupakan
pembatasan % \index{restriction}\index{map!restriction}\appendrefs{map restrictions}\spacefactor=1000 %
pemetaan pertama pada $\genrangespace{t}$.
\end{remark}

Bagi pangkat antara $j=0$ dan~$j=n$, ruang $V$ mungkin bukan jumlah langsung
$\rangespace{t^j}$ dan $\nullspace{t^j}$.
Contoh berikut menunjukkan bahwa keduanya dapat memiliki irisan nontrivial.

\begin{example}   \label{FirstNilMap}
Tinjau transformasi pada \( \C^2 \) yang didefinisikan oleh
aksi berikut pada anggota-anggota basis standar.
\begin{equation*}
  \colvec[r]{1 \\ 0}
    \mapsunder{n}
    \colvec[r]{0 \\ 1}
  \quad
  \colvec[r]{0 \\ 1}
    \mapsunder{n}
    \colvec[r]{0 \\ 0}
  \qquad
  N=\rep{n}{\stdbasis_2,\stdbasis_2}=\begin{mat}[r]
    0  &0  \\
    1  &0
  \end{mat}
\end{equation*}
Ini adalah \definend{pemetaan geser}\index{geser}\index{transformasi!geser}
karena pemetaan tersebut menggeser entri-entri ke bawah, dan entri paling bawah
bergeser sepenuhnya keluar dari vektor.  
\begin{equation*}
  \colvec{x \\ y}\mapsto \colvec{0 \\ x}
\end{equation*}
Pada basis tersebut, aksi pemetaan ini menghasilkan suatu
\definend{untaian}.\index{untaian!vektor-vektor basis}
\begin{equation*}
  \begin{strings}{ccccc}
     \colvec{1 \\ 0} &\mapsto &\colvec{0 \\ 1} &\mapsto &\colvec{0 \\ 0}
  \end{strings}
  \qquad\text{yakni}\qquad
  \begin{strings}{ccccc}
     \vec{e}_1 &\mapsto &\vec{e}_2 &\mapsto &\zero
  \end{strings}
\end{equation*}
Pemetaan ini memberikan cara alami untuk memperoleh suatu vektor yang berada
sekaligus dalam ruang hasil dan ruang nol; gambaran untaian tersebut menunjukkan
salah satu vektor demikian.
\begin{equation*}
  \vec{e}_2=\colvec[r]{0 \\ 1}
\end{equation*}
Perhatikan pula bahwa walaupun $n$ bukan pemetaan nol,
fungsi $n^2=\composed{n}{n}$ merupakan pemetaan nol.
\end{example}

\begin{example}  \label{NilIndexFourOnCFour}
Fungsi linear \( \map{\hat{n}}{\C^4}{\C^4} \) yang aksinya pada
\( \stdbasis_4 \) diberikan oleh untaian
\begin{equation*}
  \begin{strings}{ccccccccc}
     \vec{e}_1 &\mapsto &\vec{e}_2
          &\mapsto &\vec{e}_3
          &\mapsto &\vec{e}_4
          &\mapsto &\zero
  \end{strings}
\end{equation*}
memiliki
\( \rangespace{\hat{n}}\intersection\nullspace{\hat{n}} \) yang sama dengan
rentang \( \spanof{\set{\vec{e}_4}} \),
memiliki \( \rangespace{\hat{n}^2}\intersection\nullspace{\hat{n}^2}=
  \spanof{\set{\vec{e}_3,\vec{e}_4}} \),
dan memiliki \( \rangespace{\hat{n}^3}\intersection\nullspace{\hat{n}^3}=
    \spanof{\set{\vec{e}_4}} \).
Representasi matriksnya seluruhnya nol, kecuali beberapa entri satu pada
subdiagonal.
\begin{equation*}
  \hat{N}=\rep{\hat{n}}{\stdbasis_4,\stdbasis_4}
  =\begin{mat}[r]
    0  &0  &0  &0 \\
    1  &0  &0  &0 \\
    0  &1  &0  &0 \\
    0  &0  &1  &0 \\
  \end{mat}
\end{equation*}
Walaupun $\hat{n}$ bukan pemetaan nol, demikian pula $\hat{n}^2$ dan
$\hat{n}^3$, fungsi $\hat{n}^4$ merupakan fungsi nol.
\end{example}

\begin{example} \label{ThirdNilMap}
Transformasi dapat bertindak melalui lebih dari satu untaian.
Transformasi \( t \) bertindak pada basis
\( B=\sequence{\vec{\beta}_1,\dots,\vec{\beta}_5} \) menurut
\begin{equation*}
   \begin{strings}{ccccccc}
    \vec{\beta}_1 &\mapsto &\vec{\beta}_2 &\mapsto &\vec{\beta}_3
        &\mapsto &\zero \\
    \vec{\beta}_4 &\mapsto &\vec{\beta}_5 &\mapsto &\zero
  \end{strings}
\end{equation*}
dan, misalnya, memiliki $\vec{\beta}_3$ dalam irisan ruang hasil dan ruang
nolnya.
Untaian-untaian tersebut memperjelas bahwa $t^3$ merupakan pemetaan nol.
Pemetaan ini direpresentasikan oleh matriks yang seluruhnya nol, kecuali
blok-blok entri satu pada subdiagonal,
\begin{equation*}
  \rep{t}{B,B}=
  \begin{pmat}{rrr|rr}
     0  &0  &0  &0  &0  \\
     1  &0  &0  &0  &0  \\
     0  &1  &0  &0  &0  \\ \hline
     0  &0  &0  &0  &0  \\
     0  &0  &0  &1  &0
  \end{pmat}
\end{equation*}
(garis-garisnya hanya mengatur blok-blok tersebut secara visual).
\end{example}

Dalam contoh-contoh tersebut, semua vektor pada akhirnya ditransformasikan
menjadi nol.

\begin{definition} \label{def:nilpotent} \index{nilpoten!definisi}
Transformasi \definend{nilpoten}\index{transformasi!nilpoten}%
\index{nilpoten!transformasi}
adalah transformasi yang suatu pangkatnya merupakan pemetaan nol.
\definend{Matriks nilpoten}\index{matriks!nilpoten}%
\index{nilpoten!matriks}
adalah matriks yang suatu pangkatnya merupakan matriks nol.
Dalam kedua kasus, pangkat terkecil demikian disebut
\definend{indeks nilpotensi}.%
\index{nilpotensi!indeks}\index{indeks nilpotensi}
\end{definition}

\begin{example}
Dalam \nearbyexample{FirstNilMap}, indeks nilpotensinya adalah dua.
Dalam \nearbyexample{NilIndexFourOnCFour}, indeksnya empat.
Dalam \nearbyexample{ThirdNilMap}, indeksnya tiga.
\end{example}

\begin{example}
Pemetaan diferensiasi \( \map{d/dx}{\polyspace_2}{\polyspace_2} \)
nilpoten dengan indeks tiga karena turunan ketiga sebarang polinomial
kuadrat adalah nol.
Aksi pemetaan ini dideskripsikan oleh untaian
$x^2\mapsto 2x\mapsto 2\mapsto 0$
dan pemilihan basis \( B=\sequence{x^2,2x,2} \)
memberikan representasi berikut.
\begin{equation*}
  \rep{d/dx}{B,B}=
  \begin{mat}[r]
     0  &0  &0  \\
     1  &0  &0  \\
     0  &1  &0
  \end{mat}
\end{equation*}
\end{example}

Tidak semua matriks nilpoten seluruhnya nol kecuali blok-blok entri satu pada
subdiagonal.

\begin{example} \label{ex:NilMatNotCanon}
Untuk matriks $\hat{N}$ dari \nearbyexample{NilIndexFourOnCFour}
dan basis empat vektor berikut,
\begin{equation*}
  D=\sequence{\colvec[r]{1 \\ 0 \\ 1 \\ 0},
              \colvec[r]{0 \\ 2 \\ 1 \\ 0},
              \colvec[r]{1 \\ 1 \\ 1 \\ 0},
              \colvec[r]{0 \\ 0 \\ 0 \\ 1}}
\end{equation*}
operasi perubahan basis menghasilkan representasi berikut terhadap \( D,D \).
\begin{equation*}
  \begin{mat}[r]
    1  &0  &1 &0 \\
    0  &2  &1 &0 \\
    1  &1  &1 &0 \\
    0  &0  &0 &1
  \end{mat}
  \begin{mat}[r]
    0  &0  &0 &0 \\
    1  &0  &0 &0 \\
    0  &1  &0 &0 \\
    0  &0  &1 &0
  \end{mat}
  \begin{mat}[r]
    1  &0  &1 &0 \\
    0  &2  &1 &0 \\
    1  &1  &1 &0 \\
    0  &0  &0 &1
  \end{mat}^{-1}\!\!
  =
  \begin{mat}[r]
   -1  &0  &1   &0 \\
   -3  &-2 &5   &0 \\
   -2  &-1  &3  &0 \\
    2  &1   &-2 &0
  \end{mat}
\end{equation*}
Matriks baru tersebut nilpoten; pangkat keempatnya adalah matriks nol.
Kita dapat memverifikasinya dengan perhitungan yang melelahkan, atau cukup
mengamati bahwa matriks itu nilpoten karena pangkat keempatnya serupa dengan
\( \hat{N}^4 \), yaitu matriks nol, dan satu-satunya matriks yang serupa dengan
matriks nol adalah matriks nol itu sendiri.
\begin{equation*}
   (P\hat{N}P^{-1})^4
   =P\hat{N}P^{-1}\cdot P\hat{N}P^{-1}\cdot P\hat{N}P^{-1}\cdot P\hat{N}P^{-1}
   =P\hat{N}^4P^{-1}
\end{equation*}
\end{example}

Tujuan subbagian ini adalah menunjukkan bahwa contoh sebelumnya bersifat
prototipikal: setiap matriks nilpoten serupa dengan suatu matriks yang seluruhnya
nol, kecuali blok-blok entri satu pada subdiagonal.

\begin{definition}
Misalkan \( t \) suatu transformasi nilpoten pada \( V \).
Suatu \definend{untaian \( t \) yang dihasilkan oleh
\( \vec{v}\in V \)}\index{untaian}
adalah barisan
\( \sequence{\vec{v},t(\vec{v}),\ldots,t^{k-1}(\vec{v})} \)
sedemikian sehingga $t^k(\vec{v})=\zero$.
Suatu \definend{basis untaian \( t \)}\index{basis!untaian}\index{untaian!basis}
adalah basis yang merupakan konkatenasi untaian-untaian \( t \).
\end{definition}

\noindent (Untaian-untaian tidak dapat membentuk basis melalui konkatenasi
kecuali jika saling lepas, karena sebuah basis tidak dapat memuat vektor yang berulang.)

\begin{example}
Pemetaan linear \( \map{t}{\C^3}{\C^3} \)
\begin{equation*}
  \colvec{x \\ y \\ z}\mapsunder{t}\colvec{y \\ z \\ 0}
\end{equation*}
bersifat nilpoten dengan indeks~$3$.
\begin{equation*}
  \colvec{x \\ y \\ z}
  \mapsunder{t}\colvec{y \\ z \\ 0}
  \mapsunder{t}\colvec{z \\ 0 \\ 0}
  \mapsunder{t}\colvec{0 \\ 0 \\ 0}
\end{equation*}
Berikut ini adalah suatu untaian~$t$.
\begin{equation*}
  \sequence{\colvec{0 \\ 0 \\ 1},
            \colvec{0 \\ 1 \\ 0},
            \colvec{1 \\ 0 \\ 0}}
\end{equation*}
Karena barisan tersebut merupakan basis,
barisan itu adalah basis untaian~$t$ bagi ruang~$\C^3$.
\end{example}

Jika barisan pada contoh sebelumnya adalah
$\sequence{\vec{\beta}_1,
             \vec{\beta}_2,
             \vec{\beta}_3}$,
untaian~$t$ lainnya adalah
$\sequence{\vec{\beta}_2,
             \vec{\beta}_3}$.
Namun, tentu saja barisan kedua bukan basis bagi~$\C^3$.
Dalam pengertian ini, untaian-untaian~$t$ dalam suatu basis untaian~$t$
harus maksimal.

\begin{example}
Pemetaan linear diferensiasi
\( \map{d/dx}{\polyspace_2}{\polyspace_2} \)
bersifat nilpoten.
Barisan
\( \sequence{x^2, 2x, 2} \)
adalah untaian \( d/dx \) dengan panjang~\( 3 \);
khususnya, untaian ini memenuhi syarat $d/dx(2)=0$.
Karena merupakan basis, barisan tersebut
adalah basis untaian \( d/dx \) bagi \( \polyspace_2\).
\end{example}

\begin{example}
Dalam \nearbyexample{ThirdNilMap}, kita dapat mengonkatkan untaian-untaian~$t$
$\sequence{\vec{\beta}_1,\vec{\beta}_2,\vec{\beta}_3}$ dan
$\sequence{\vec{\beta}_4,\vec{\beta}_5}$
untuk membentuk basis bagi domain~$t$.
\end{example}

\begin{lemma}  \label{le:LongestTowerIsIndex}
Jika suatu ruang mempunyai basis untaian \( t \), maka indeks nilpotensi~$t$
sama dengan panjang untaian terpanjang dalam basis tersebut.
\end{lemma}

\begin{proof}
Misalkan ruang tersebut mempunyai basis untaian~$t$ dan indeks
nilpotensi~$t$ adalah~$k$.
Maka \( t^k \) memetakan setiap vektor ke \( \zero \).
Hal ini berlaku pula bagi
vektor awal setiap untaian,
sehingga setiap untaian dalam basis untaian itu panjangnya paling besar~$k$.

Sebaliknya, andaikan ruang tersebut mempunyai basis untaian~$t$, yaitu~$B$,
yang semua untaiannya lebih pendek daripada~\( k \).
Karena indeks nilpotensi~$t$ adalah~$k$, terdapat \( \vec{v} \)
sedemikian sehingga \( t^{k-1}(\vec{v})\neq\zero \).
Nyatakan $\vec{v}$ sebagai kombinasi linear unsur-unsur~$B$
lalu terapkan \( t^{k-1} \).
Menurut pengandaian kita, \( t^{k-1} \) memetakan setiap unsur~$B$ ke \( \zero \).
Karena itu, pemetaan tersebut membawa setiap suku dalam kombinasi linear ke~$\zero$,
bertentangan dengan fakta bahwa \( \vec{v} \) tidak dipetakan ke \( \zero \).
\end{proof}

Kita akan menunjukkan bahwa setiap
pemetaan nilpoten mempunyai basis untaian yang terkait, yakni basis yang
terdiri atas untaian-untaian saling lepas.
% Then our goal theorem, that every
% nilpotent matrix is similar to one that is
% all zeros except for blocks of subdiagonal ones, is immediate,
 %as in \nearbyexample{ThirdNilMap}.

Untuk memahami gagasan utama argumen ini, bayangkan bahwa
kita ingin membuat suatu contoh tandingan, yakni pemetaan yang nilpoten tetapi
tidak mempunyai basis untaian saling lepas yang terkait.
Kita mungkin mencoba membuat pemetaan seperti
\( \map{t}{\C^5}{\C^5} \) dengan aksi berikut.
\begin{center}
  \begin{minipage}{1in}
  \setlength{\unitlength}{1pt}
  $\begin{strings}{ccccc}
     \begin{picture}(4,15)
       \put(0,14){$\vec{e}_1$}
       \put(0,-12){$\vec{e}_2$}
     \end{picture}
     &\begin{picture}(8,10)
       \put(0,8){\rotatebox{-30}{$\mapsto$}}%  
       \put(0,-8){\rotatebox{30}{$\mapsto$}}
     \end{picture}
     &\vec{e}_3 &\mapsto &\zero  \\[4ex]
     \vec{e}_4 &\mapsto &\vec{e}_5 &\mapsto &\zero  
  \end{strings}$
  \end{minipage}
  \hspace*{3em}
  $\rep{t}{\stdbasis_5,\stdbasis_5}=
  \begin{mat}[r]
    0  &0  &0  &0  &0  \\
    0  &0  &0  &0  &0  \\
    1  &1  &0  &0  &0  \\
    0  &0  &0  &0  &0  \\
    0  &0  &0  &1  &0
  \end{mat}$
\end{center}
Namun, fakta bahwa basis yang ditampilkan tidak saling lepas tidak berarti
bahwa tidak ada basis lain yang terdiri atas untaian-untaian saling lepas.

Untuk menghasilkan basis semacam itu bagi pemetaan ini,
mula-mula kita tentukan banyak dan panjang untaiannya.
Perhatikan bahwa indeks nilpotensi~$t$ adalah dua.
\nearbylemma{le:LongestTowerIsIndex} menyatakan bahwa dalam suatu basis
untaian saling lepas, sedikitnya satu untaian mempunyai panjang dua.
Ada lima unsur basis, sehingga jika terdapat basis untaian saling lepas,
pemetaan tersebut harus beraksi dengan salah satu dari cara berikut.
\begin{equation*}
  \begin{strings}{ccccc}
    \vec{\beta}_1 &\mapsto &\vec{\beta}_2 &\mapsto &\zero  \\
    \vec{\beta}_3 &\mapsto &\vec{\beta}_4 &\mapsto &\zero  \\
    \vec{\beta}_5 &\mapsto &\zero
  \end{strings}
  \hspace*{3em}
  \begin{strings}{ccccc}
    \vec{\beta}_1 &\mapsto &\vec{\beta}_2 &\mapsto &\zero  \\
    \vec{\beta}_3 &\mapsto &\zero   \\
    \vec{\beta}_4 &\mapsto &\zero   \\
    \vec{\beta}_5 &\mapsto &\zero
  \end{strings}
\end{equation*}
Sekarang, inilah titik kuncinya.
Transformasi dengan aksi di sebelah kiri mempunyai
ruang nol berdimensi tiga, karena sebanyak itulah vektor basis yang
dipetakan ke nol.
Transformasi dengan aksi di sebelah kanan mempunyai ruang nol
berdimensi empat.
Dengan representasi matriks di atas, kita dapat menentukan bentuk mana
di antara dua kemungkinan tersebut yang benar.
\begin{equation*}
  \nullspace{t}=
  \set{\colvec{x \\ -x \\ z \\ 0 \\ r}\suchthat x,z,r\in\C }
\end{equation*}
Ruang ini berdimensi tiga,
yang berarti bahwa dari dua bentuk basis untaian saling lepas di atas,
basis bagi \( t \) mempunyai bentuk di sebelah kiri.

Untuk menghasilkan basis untaian bagi~$t$, mula-mula
pilih \( \vec{\beta}_2 \) dan \( \vec{\beta}_4 \) dari
\( \rangespace{t}\intersection\nullspace{t} \).
\begin{equation*}
  \vec{\beta}_2=\colvec[r]{0 \\ 0 \\ 1 \\ 0 \\ 0}\qquad
  \vec{\beta}_4=\colvec[r]{0 \\ 0 \\ 0 \\ 0 \\ 1}
\end{equation*}
(Pilihan lain juga dimungkinkan; pastikan saja bahwa himpunan
\( \set{\vec{\beta}_2,\vec{\beta}_4} \) bebas linear.)
Untuk \( \vec{\beta}_5 \), pilih suatu vektor dari \( \nullspace{t} \)
yang tidak berada dalam rentang \( \set{ \vec{\beta}_2,\vec{\beta}_4 } \).
\begin{equation*}
  \vec{\beta}_5=\colvec[r]{1 \\ -1 \\ 0 \\ 0 \\ 0}
\end{equation*}
Terakhir, ambil \( \vec{\beta}_1 \) dan \( \vec{\beta}_3 \) sedemikian sehingga
\( t(\vec{\beta}_1)=\vec{\beta}_2 \) dan
\( t(\vec{\beta}_3)=\vec{\beta}_4 \).
\begin{equation*}
  \vec{\beta}_1=\colvec[r]{0 \\ 1 \\ 0 \\ 0 \\ 0}\qquad
  \vec{\beta}_3=\colvec[r]{0 \\ 0 \\ 0 \\ 1 \\ 0}
\end{equation*}
Dengan demikian, kita memperoleh basis untaian
\( B=\sequence{\vec{\beta}_1,\ldots,\vec{\beta}_5} \)
dan, terhadap basis tersebut,
matriks~$t$ mempunyai blok-blok entri~$1$ pada subdiagonal.
\begin{equation*}
  \rep{t}{B,B}=
  \begin{pmat}{rr|rr|r}
    0  &0  &0  &0  &0  \\
    1  &0  &0  &0  &0  \\  \hline
    0  &0  &0  &0  &0  \\
    0  &0  &1  &0  &0  \\  \hline
    0  &0  &0  &0  &0
  \end{pmat}
\end{equation*}

\begin{theorem}
\label{th:NilMapHasStrBas}
Setiap transformasi nilpoten~$t$ mempunyai suatu basis untaian \( t \) yang terkait.
Meskipun basisnya tidak unik, banyak
dan panjang untaiannya ditentukan oleh \( t \).
\end{theorem}

Gambar berikut mengilustrasikan bukti, yang menjelaskan tiga jenis
vektor basis.
Vektor tersebut berada dalam persegi jika merupakan anggota
ruang nol, dan dalam lingkaran jika bukan.
\begin{equation*}
   \begin{strings}{ccccccccccccccccccc}
     \digitincirc{3}
         &\mapsto &\digitincirc{1} &\mapsto
         &\cdots &  &  &  &  &  &  &  &\cdots
         &\mapsto &\digitincirc{1}
         &\mapsto &\digitinsq{1} &\mapsto &\zero  \\[.75ex]
     \digitincirc{3}
         &\mapsto &\digitincirc{1} &\mapsto
         &\cdots &  &  &  &  &  &  &  &\cdots
         &\mapsto &\digitincirc{1}
         &\mapsto &\digitinsq{1} &\mapsto &\zero  \\[.75ex]
         &\smash{\vdotswithin{\mapsto}}  \\
     \digitincirc{3}
         &\mapsto &\digitincirc{1} &\mapsto
         &\cdots &
         &\mapsto &\digitincirc{1}
         &\mapsto &\digitinsq{1} &\mapsto &\zero  \\[1ex]
     \digitinsq{2} &\mapsto &\zero \\[.75ex]
         &\smash{\vdotswithin{\mapsto}}  \\
     \digitinsq{2} &\mapsto &\zero
   \end{strings}
\end{equation*}

\begin{proof}
Tetapkan suatu ruang vektor~$V$.
Kita akan berargumen dengan induksi pada indeks nilpotensi.
Jika pemetaan $\map{t}{V}{V}$
mempunyai indeks nilpotensi~\( 1 \), maka pemetaan itu adalah pemetaan nol
dan sebarang basis merupakan basis untaian,
$\vec{\beta}_1\mapsto\zero$, \ldots,
$\vec{\beta}_n\mapsto\zero$.

Untuk langkah induksi, andaikan teorema berlaku bagi setiap transformasi
$\map{t}{V}{V}$
dengan indeks nilpotensi dari~$1$ sampai dengan~\( k-1 \) (dengan $k>1$),
dan tinjau kasus berindeks~$k$.

Pengamatan yang menjalankan induksi adalah bahwa
pembatasan~$t$ pada ruang hasil~\( \rangespace{t} \)
juga nilpoten, dengan indeks \( k-1 \).
Karena itu, terapkan hipotesis induksi untuk memperoleh basis untaian bagi
\( \rangespace{t} \),
yang banyak dan panjang untaiannya ditentukan oleh \( t \).
\begin{equation*}
  B=\cat{\cat{\sequence{\vec{\beta}_1,t(\vec{\beta}_1),\dots,
     t^{h_1}(\vec{\beta}_1)}}{
  \cat{\sequence{\vec{\beta}_2,\ldots,t^{h_2}(\vec{\beta}_2)}}}{\cdots}}{
  \sequence{\vec{\beta}_i,\ldots,t^{h_i}(\vec{\beta}_i)} }
\end{equation*}
Kita memakai~$i$ untuk menyatakan banyak untaian.
Pada gambar di atas, vektor-vektor ini adalah vektor jenis~\( 1 \).

Mengambil vektor terakhir dari setiap untaian
menghasilkan basis
\( C=\sequence{t^{h_1}(\vec{\beta}_1),\dots,t^{h_i}(\vec{\beta}_i)} \)
bagi irisan \( \rangespace{t}\intersection\nullspace{t} \).
Hal ini karena suatu anggota \( \rangespace{t} \) dipetakan ke nol jika dan
hanya jika ia merupakan kombinasi linear vektor-vektor basis yang dipetakan
ke nol.
Pada gambar, vektor-vektor ini ditunjukkan sebagai \( 1 \) di dalam persegi.

Sekarang perluas \( C \) menjadi basis bagi seluruh ruang nol, \( \nullspace{t} \).
\begin{equation*}
  \hat{C}=\cat{C}{\sequence{\vec{\xi}_1,\dots,\vec{\xi}_p}}
\end{equation*}
Meskipun vektor-vektor \( \vec{\xi} \) tidak ditentukan secara unik oleh~$t$,
banyaknya ditentukan secara unik, yakni dimensi
\( \nullspace{t} \) dikurangi dimensi
\( \rangespace{t}\intersection\nullspace{t} \).
Pada gambar, vektor-vektor~$\vec{\xi}$ adalah vektor jenis~\( 2 \),
sehingga \( \hat{C} \) adalah himpunan vektor dalam persegi.

Terakhir, \( \cat{B}{\sequence{\vec{\xi}_1,\dots,\vec{\xi}_p}} \)
adalah basis bagi \( \rangespace{t}+\nullspace{t} \).
Hal ini karena jumlah suatu unsur ruang hasil dengan suatu unsur ruang nol
dapat direpresentasikan menggunakan unsur-unsur \( B \) untuk bagian ruang
hasil, bersama vektor-vektor~$\xi$ untuk setiap bagian yang berasal dari
\( \nullspace{t}-\rangespace{t} \).
Perhatikan bahwa
\begin{align*}
  \dim\big(\rangespace{t}+\nullspace{t}\big)
  &=
  \dim (\rangespace{t})+\dim (\nullspace{t})
    -\dim\big(\rangespace{t}\intersection\nullspace{t}\big)  \\
  &=
  \rank (t)+\nullity (t)-i          \\
  &=
  \dim (V)-i
\end{align*}
sehingga kita dapat memperluas basis
\( \cat{B}{\sequence{\vec{\xi}_1,\dots,\vec{\xi}_p}} \) 
menjadi basis untaian~$t$ bagi seluruh
\( V \) dengan menambahkan~\( i \) vektor lagi.
Untuk menghasilkannya,
ingat bahwa setiap \( \vec{\beta}_1,\dots,\vec{\beta}_i \) berada
di ruang hasil, \( \rangespace{t} \),
sehingga gunakan vektor-vektor
\( \vec{v}_1,\dots,\vec{v}_i\in V \) sedemikian sehingga
\( t(\vec{v}_1)=\vec{\beta}_1,\dots,t(\vec{v}_i)=\vec{\beta}_i \).
Pemeriksaan bahwa perluasan ini mempertahankan kebebasan linear terdapat dalam
\nearbyexercise{exer:NilMapWillHaveStrBas}.
Pada gambar, \( \vec{v}_1,\dots,\vec{v}_i \) adalah vektor-vektor \( 3 \).
\end{proof}

\begin{corollary} \label{cor:NilpotentMatCanonForm}
\index{nilpoten!bentuk kanonik untuk}
\index{transformasi!nilpoten!wakil kanonik}
\index{bentuk kanonik!untuk matriks nilpoten}
Setiap matriks nilpoten serupa dengan matriks yang seluruh entrinya nol,
kecuali blok-blok entri satu pada subdiagonal.
Artinya, setiap pemetaan nilpoten direpresentasikan terhadap suatu basis
oleh matriks semacam itu.
\end{corollary}

Bentuk ini unik dalam arti bahwa jika suatu matriks nilpoten serupa dengan dua
matriks semacam itu, maka kedua matriks tersebut hanya berbeda dalam urutan bloknya.
Jadi, bentuk ini merupakan
bentuk kanonik bagi kelas-kelas keserupaan matriks nilpoten,
asalkan kita mengurutkan blok-bloknya, misalnya dari yang terpanjang ke yang terpendek.

\begin{example}
Matriks
\begin{equation*}
  M=\begin{mat}[r]
      1  &-1  \\
      1  &-1
    \end{mat}
\end{equation*}
mempunyai indeks nilpotensi dua, sebagaimana ditunjukkan oleh perhitungan berikut.
\begin{center}
  \begin{tabular}{c|cc}
    \multicolumn{1}{c}{\textit{pangkat} \( p \)}  &\( M^p \)  &\( \nullspace{M^p}  \)   \\  
    \hline
    \( 1 \)
    &\(  
       M=\matrixvenlarge{\begin{mat}[r]
         1  &-1  \\
         1  &-1
       \end{mat}}  \)
    &\( \set{\matrixvenlarge{\colvec{x \\ x}}\suchthat
                               x\in\C}  \)   \\
    \( 2 \)
    &\(  M^2=\matrixvenlarge{\begin{mat}[r]
         0  &0   \\
         0  &0
       \end{mat}}  \)
    &\( \C^2  \)
  \end{tabular}
\end{center}
Karena matriks itu berukuran $\nbyn{2}$, setiap transformasi yang
direpresentasikannya bekerja pada ruang berdimensi dua.
Ruang nol setelah pemetaan diterapkan satu kali,
$\nullspace{m}$, berdimensi satu, sedangkan
ruang nol setelah diterapkan dua kali, $\nullspace{m^2}$, berdimensi dua.
Jadi, aksi $m$ pada suatu basis untaian adalah
$\vec{\beta}_1\mapsto\vec{\beta}_2\mapsto\zero$ 
dan bentuk kanonik matriksnya adalah berikut ini.
\begin{equation*}
  N=\begin{mat}[r]
    0  &0  \\
    1  &0
  \end{mat}
\end{equation*}

Kita dapat memperlihatkan basis untaian semacam itu,
beserta matriks-matriks perubahan basis yang membuktikan keserupaan matriks
antara $M$ dan~$N$.
Andaikan \( M \) merepresentasikan $\map{m}{\C^2}{\C^2}$
terhadap basis standar.
(Kita dapat saja memilih $M$ sebagai representasi terhadap basis lain,
tetapi basis standar lebih mudah digunakan.)
Pilih \( \vec{\beta}_2\in\nullspace{m} \).
Pilih pula \( \vec{\beta}_1 \)
sedemikian sehingga \( m(\vec{\beta}_1)=\vec{\beta}_2 \).
\begin{equation*}
  \vec{\beta}_2=\colvec[r]{1 \\ 1}
  \qquad
  \vec{\beta}_1=\colvec[r]{1 \\ 0}
\end{equation*}
Untuk matriks perubahan basis, ingat kembali diagram keserupaan berikut.
\begin{equation*}
  \begin{CD}
    \C^2_{\wrt{\stdbasis_2}}      @>m>M>        \C^2_{\wrt{\stdbasis_2}}     \\
    @V\scriptstyle\identity V\scriptstyle PV  @V\scriptstyle\identity V\scriptstyle PV \\
    \C^2_{\wrt{B}}                 @>m>N>         \C^2_{\wrt{B}}
  \end{CD}
\end{equation*}
Bentuk kanoniknya adalah \( \rep{m}{B,B}=PMP^{-1} \), dengan
\begin{equation*}
   P^{-1} %=\bigl(\rep{\identity}{\stdbasis_2,B}\bigr)^{-1}
         =\rep{\identity}{B,\stdbasis_2}
         =\begin{mat}[r]
            1  &1  \\
            0  &1
          \end{mat}
   \qquad
   P=(P^{-1})^{-1}
         =\begin{mat}[r]
            1  &-1  \\
            0  &1
          \end{mat}
\end{equation*}
dan verifikasi perhitungan matriksnya bersifat rutin.
\begin{equation*}
  \begin{mat}[r]
    1  &-1  \\
    0  &1
  \end{mat}
  \begin{mat}[r]
    1  &-1  \\
    1  &-1
  \end{mat}
  \begin{mat}[r]
    1  &1  \\
    0  &1
  \end{mat}=
  \begin{mat}[r]
    0  &0  \\
    1  &0
  \end{mat}
\end{equation*}
%which does indeed describe the effect of $m$ on the string basis.
\end{example}

\begin{example} \label{ex:FiveByNilStrBas}
Matriks berikut
\begin{equation*}
  \begin{mat}[r]
     0  &0  &0  &0  &0  \\
     1  &0  &0  &0  &0  \\
     -1 &1  &1  &-1 &1  \\
     0  &1  &0  &0  &0  \\
     1  &0  &-1 &1  &-1
  \end{mat}
\end{equation*}
bersifat nilpoten dengan indeks~$3$.
% \newlength{\extramatrixvspace}
% \setlength{\extramatrixvspace}{.75ex}
% \newcommand{\matrixvenlarge}[1]{\raisebox{-0.5\height}{\vbox{
%        \vspace*{\extramatrixvspace}
%        \hbox{$#1$}
%         \vspace*{\extramatrixvspace}
%         }}}
\begin{center} 
  \begin{tabular}{c|cc}
    \multicolumn{1}{c}{\textit{pangkat} \( p \)}  &\( N^p \)  &\( \nullspace{N^p}  \)   \\  
    \hline
    \( 1 \)
    &\matrixvenlarge{
         \begin{mat}[r]
         0  &0  &0  &0  &0  \\
         1  &0  &0  &0  &0  \\
         -1 &1  &1  &-1 &1  \\
         0  &1  &0  &0  &0  \\
         1  &0  &-1 &1  &-1
         \end{mat}}
    &\( 
        \set{\matrixvenlarge{\colvec{0 \\ 0 \\ u-v \\ u \\ v}} 
           \suchthat u,v\in\C}  \)                      \\
    \( 2 \)
    &\(
       \matrixvenlarge{\begin{mat}[r]
         0  &0  &0  &0  &0  \\
         0  &0  &0  &0  &0  \\
         1  &0  &0  &0  &0  \\
         1  &0  &0  &0  &0  \\
         0  &0  &0  &0  &0
       \end{mat}}  \)
    &\( 
        \set{\matrixvenlarge{\colvec{0 \\ y \\ z \\ u \\ v}}
                       \suchthat y,z,u,v\in\C}  \)  \\
    \( 3 \)
    &\textit{--matriks nol--}
    &\( \C^5 \)
  \end{tabular}
\end{center}
Tabel tersebut memberi tahu kita hal berikut tentang sebarang basis untaian:~ruang
nol setelah satu kali penerapan pemetaan berdimensi dua, sehingga dua vektor
basis dipetakan langsung ke nol; ruang nol setelah penerapan kedua berdimensi
empat, sehingga dua vektor basis tambahan dipetakan ke nol pada iterasi kedua;
dan ruang nol setelah tiga kali penerapan berdimensi lima, sehingga satu vektor
basis yang tersisa dipetakan ke nol dalam tiga langkah.
\begin{equation*}
  \begin{strings}{ccccccc}
    \vec{\beta}_1 &\mapsto &\vec{\beta}_2 &\mapsto &\vec{\beta}_3
       &\mapsto &\zero  \\
    \vec{\beta}_4 &\mapsto &\vec{\beta}_5 &\mapsto &\zero
  \end{strings}
\end{equation*}
Untuk menghasilkan basis semacam itu, mula-mula pilih dua vektor dari
\( \nullspace{n} \) yang membentuk himpunan bebas linear.
\begin{equation*}
   \vec{\beta}_3=\colvec[r]{0 \\ 0 \\ 1 \\ 1 \\ 0} \quad
   \vec{\beta}_5=\colvec[r]{0 \\ 0 \\ 0 \\ 1 \\ 1}
\end{equation*}
Kemudian tambahkan \( \vec{\beta}_2,\vec{\beta}_4\in\nullspace{n^2} \)
sedemikian sehingga \( n(\vec{\beta}_2)=\vec{\beta}_3 \) dan
\( n(\vec{\beta}_4)=\vec{\beta}_5 \).
\begin{equation*}
   \vec{\beta}_2=\colvec[r]{0 \\ 1 \\ 0 \\ 0 \\ 0} \quad
   \vec{\beta}_4=\colvec[r]{0 \\ 1 \\ 0 \\ 1 \\ 0}
\end{equation*}
Akhiri dengan menambahkan \( \vec{\beta}_1 \)
sedemikian sehingga \( n(\vec{\beta}_1)=\vec{\beta}_2 \).
\begin{equation*}
   \vec{\beta}_1=\colvec[r]{1 \\ 0 \\ 1 \\ 0 \\ 0}
\end{equation*}
\end{example}






\begin{exercises}
   \recommended \item \label{exer:IndNilLftShift}
     Berapakah indeks nilpotensi operator \definend{geser ke kanan},
     yang di sini bekerja pada ruang tripel bilangan real?
      \begin{equation*}
         (x,y,z)\mapsto(0,x,y)
      \end{equation*}
      \begin{answer}
        Tiga.
        Indeksnya sedikitnya tiga karena $\ell^2(\,(1,1,1)\,)=(0,0,1)\neq \zero$.
        Indeksnya paling besar tiga karena
        $(x,y,z)\mapsto (0,x,y)\mapsto (0,0,x)\mapsto (0,0,0)$.
      \end{answer}
  \recommended \item 
    Untuk setiap basis untaian, nyatakan indeks nilpotensinya dan
    berikan dimensi ruang hasil serta
    ruang nol pada setiap iterasi pemetaan nilpoten tersebut.
    \begin{exparts}
      \partsitem $
        \begin{strings}{ccccccc}
           \vec{\beta}_1 &\mapsto &\vec{\beta}_2 &\mapsto &\zero  \\
           \vec{\beta}_3 &\mapsto &\vec{\beta}_4 &\mapsto &\zero  
         \end{strings}$
      \partsitem $
        \begin{strings}{ccccccc}
           \vec{\beta}_1 &\mapsto &\vec{\beta}_2 &\mapsto &\vec{\beta}_3
                &\mapsto &\zero  \\
           \vec{\beta}_4 &\mapsto &\zero \\
           \vec{\beta}_5 &\mapsto &\zero \\
           \vec{\beta}_6 &\mapsto &\zero
         \end{strings}$
      \partsitem $
        \begin{strings}{ccccccccc}
           \vec{\beta}_1 &\mapsto &\vec{\beta}_2 &\mapsto &\vec{\beta}_3
                &\mapsto &\zero  
         \end{strings}$
    \end{exparts}
    Berikan pula bentuk kanonik matriksnya.
    \begin{answer}
      \begin{exparts}
        \partsitem Domainnya berdimensi empat.
          Aksi pemetaan tersebut membawa setiap vektor dalam ruang
          $c_1\cdot \vec{\beta}_1+c_2\cdot \vec{\beta}_2
            +c_3\cdot \vec{\beta}_3+c_4\cdot \vec{\beta}_4$
          ke
          $c_1\cdot \vec{\beta}_2+c_2\cdot \zero
            +c_3\cdot \vec{\beta}_4+c_4\cdot \zero
           =c_1\cdot \vec{\beta}_2+c_3\cdot\vec{\beta}_4$.
          Penerapan pertama pemetaan
          membawa dua vektor basis $\vec{\beta}_2$ dan
          $\vec{\beta}_4$ ke nol,
          sehingga ruang nol berdimensi dua dan ruang hasil
          berdimensi dua.
          Pada penerapan kedua, keempat vektor basis menuju
          nol; karena itu, ruang nol dari pangkat kedua berdimensi empat,
          sedangkan ruang hasil dari pangkat kedua berdimensi nol.
          Jadi, indeks nilpotensinya adalah dua.
          Berikut bentuk kanoniknya.
          \begin{equation*}
            \begin{mat}[r]
              0  &0  &0  &0  \\
              1  &0  &0  &0  \\
              0  &0  &0  &0  \\
              0  &0  &1  &0
            \end{mat}
          \end{equation*}
        \partsitem Dimensi domain pemetaan ini adalah enam.
          Untuk pangkat pertama, dimensi ruang nol adalah empat
          dan dimensi ruang hasil adalah dua.
          Untuk pangkat kedua, dimensi ruang nol adalah lima
          dan dimensi ruang hasil adalah satu.
          Kemudian iterasi ketiga menghasilkan ruang nol berdimensi
          enam dan ruang hasil berdimensi nol.
          Indeks nilpotensinya adalah tiga, dan berikut
          bentuk kanoniknya.
          \begin{equation*}
            \begin{mat}[r]
              0  &0  &0  &0  &0  &0  \\
              1  &0  &0  &0  &0  &0  \\
              0  &1  &0  &0  &0  &0  \\
              0  &0  &0  &0  &0  &0  \\              
              0  &0  &0  &0  &0  &0  \\              
              0  &0  &0  &0  &0  &0  
            \end{mat}
          \end{equation*}
        \partsitem Dimensi domain adalah tiga, dan indeks
          nilpotensinya adalah tiga.
          Ruang nol dari pangkat pertama berdimensi satu dan ruang hasilnya
          berdimensi dua.
          Ruang nol dari pangkat kedua berdimensi dua dan ruang hasilnya
          berdimensi satu.
          Terakhir, ruang nol dari pangkat ketiga berdimensi tiga
          dan ruang hasilnya
          berdimensi nol.
          Berikut matriks bentuk kanoniknya.
          \begin{equation*}
            \begin{mat}[r]
              0  &0  &0  \\
              1  &0  &0  \\
              0  &1  &0 
            \end{mat}
          \end{equation*}
      \end{exparts}
    \end{answer}
  \item 
    Tentukan matriks-matriks berikut yang nilpoten.
    \begin{exparts*}
      \partsitem 
        $\begin{mat}[r]
           -2  &4  \\
           -1  &2
        \end{mat}$
      \partsitem 
        $\begin{mat}[r]
          3  &1  \\
          1  &3
        \end{mat}$
      \partsitem 
        $\begin{mat}[r]
          -3  &2  &1  \\
          -3  &2  &1  \\
          -3  &2  &1
        \end{mat}$
      \partsitem 
        $\begin{mat}[r]
           1  &1  &4  \\
           3  &0  &-1 \\
           5  &2  &7
        \end{mat}$
      \partsitem 
        $\begin{mat}[r]
           45  &-22  &-19  \\
           33  &-16  &-14  \\
           69  &-34  &-29
        \end{mat}$
    \end{exparts*}
    \begin{answer}
      Menurut \nearbylemma{le:RangeAndNullChains}, nulitas telah mencapai
      nilai sebesar mungkin pada iterasi ke-$n$, dengan $n$ dimensi domain.
      Jadi, untuk matriks $\nbyn{2}$,
      kita hanya perlu memeriksa apakah kuadratnya merupakan matriks nol.
      Untuk matriks $\nbyn{3}$, kita hanya perlu memeriksa kubiknya.
      \begin{exparts}
        \partsitem Ya, matriks ini nilpoten karena
          kuadratnya merupakan matriks nol.
        \partsitem Tidak, kuadratnya bukan matriks nol.
          \begin{equation*}
            \begin{mat}[r]
               3  &1  \\
               1  &3
            \end{mat}^2
            =\begin{mat}[r]
               10  &6  \\
               6   &10
            \end{mat}
          \end{equation*}
        \partsitem Ya, kubiknya merupakan matriks nol.
          Bahkan, kuadratnya nol.
        \partsitem Tidak, pangkat tiganya bukan matriks nol.
          \begin{equation*}
            \begin{mat}[r]
              1  &1  &4  \\
              3  &0  &-1 \\
              5  &2  &7
            \end{mat}^3
            =\begin{mat}[r]
              206  &86  &304  \\
               26  &8   &26   \\
              438  &180 &634
            \end{mat}
          \end{equation*}
        \partsitem Ya, kubik matriks ini merupakan matriks nol.
      \end{exparts}
      Cara lain untuk melihat bahwa matriks kedua dan keempat tidak nilpoten
      adalah dengan memperhatikan bahwa keduanya nonsingular.
    \end{answer}
  \recommended \item 
    Tentukan bentuk kanonik matriks berikut.
    \begin{equation*}
      \begin{mat}[r]
        0  &1  &1  &0  &1  \\
        0  &0  &1  &1  &1  \\
        0  &0  &0  &0  &0  \\
        0  &0  &0  &0  &0  \\
        0  &0  &0  &0  &0
      \end{mat}
    \end{equation*}
    \begin{answer} Tabel perhitungan
      \begin{center}
          \begin{tabular}{c|cc}
             \multicolumn{1}{c}{\( p \)}  &\( N^p \)  &\( \nullspace{N^p}  \) \\
             \hline
             \( 1 \)
               &\matrixvenlarge{\begin{mat}[r]
                   0  &1  &1  &0  &1  \\
                   0  &0  &1  &1  &1  \\
                   0  &0  &0  &0  &0  \\
                   0  &0  &0  &0  &0  \\
                   0  &0  &0  &0  &0
                  \end{mat}}
               &\( \set{\matrixvenlarge{\colvec{r \\ u \\ -u-v \\ u \\ v}} 
                              \suchthat r,u,v\in\C}  \) \\
             \( 2 \)
               &\matrixvenlarge{\begin{mat}[r]
                   0  &0  &1  &1  &1  \\
                   0  &0  &0  &0  &0  \\
                   0  &0  &0  &0  &0  \\
                   0  &0  &0  &0  &0  \\
                   0  &0  &0  &0  &0
                  \end{mat}}
               &\( \set{\matrixvenlarge{\colvec{r \\ s \\ -u-v \\ u \\ v}} 
                              \suchthat r,s,u,v\in\C}  \) \\
             \( 3 \)
               &\textit{--matriks nol--}
               &\( \C^5 \)
          \end{tabular}
        \end{center}
        memberikan syarat-syarat berikut bagi basis untaian:~tiga vektor basis
        dipetakan langsung ke nol, satu vektor basis lagi dipetakan ke nol
        pada penerapan kedua, dan vektor basis terakhir
        dipetakan ke nol pada penerapan ketiga.
        Jadi, basis untaiannya mempunyai bentuk berikut.
        \begin{equation*}
          \begin{strings}{ccccccc}
            \vec{\beta}_1 &\mapsto &\vec{\beta}_2 &\mapsto 
              &\vec{\beta}_3 &\mapsto &\zero               \\
            \vec{\beta}_4 &\mapsto &\zero                  \\
            \vec{\beta}_5 &\mapsto &\zero  
          \end{strings}
        \end{equation*}
        Dari sini, bentuk kanoniknya langsung diperoleh.
        \begin{equation*}
          \begin{mat}[r]
            0  &0  &0  &0  &0  \\
            1  &0  &0  &0  &0  \\
            0  &1  &0  &0  &0  \\
            0  &0  &0  &0  &0  \\
            0  &0  &0  &0  &0
          \end{mat}
        \end{equation*}
    \end{answer}
  \recommended \item 
    Tinjau matriks dari \nearbyexample{ex:FiveByNilStrBas}.
    \begin{exparts}
      \partsitem Gunakan aksi pemetaan pada basis untaian untuk
        memberikan bentuk kanoniknya.
      \partsitem Tentukan matriks-matriks perubahan basis yang membawa matriks
        tersebut ke bentuk kanonik.
      \partsitem Gunakan jawaban pada butir sebelumnya untuk memeriksa jawaban
        pada butir pertama.
    \end{exparts}
    \begin{answer}
      \begin{exparts}
        \partsitem Bentuk kanoniknya mempunyai satu blok $\nbyn{3}$ dan satu
          blok $\nbyn{2}$
          \begin{equation*}
            \begin{pmat}{rrr|rr}
              0  &0  &0  &0  &0  \\
              1  &0  &0  &0  &0  \\
              0  &1  &0  &0  &0  \\ \hline
              0  &0  &0  &0  &0  \\
              0  &0  &0  &1  &0  \\
            \end{pmat}
          \end{equation*}
          yang masing-masing bersesuaian dengan untaian berpanjang tiga dan
          untaian berpanjang dua dalam basis tersebut.
        \partsitem Andaikan $N$ adalah representasi pemetaan yang mendasarinya
          terhadap basis standar.
          Misalkan $B$ adalah basis tujuan perubahan.
          Menurut diagram keserupaan
          \begin{equation*}
            \begin{CD}
              \C^5_{\wrt{\stdbasis_5}}      
                  @>n>N>        
                  \C^5_{\wrt{\stdbasis_5}}     \\
             @V\scriptstyle\identity V\scriptstyle PV  
                 @V\scriptstyle\identity V\scriptstyle PV \\
             \C^5_{\wrt{B}}                 
                 @>n>>         
             \C^5_{\wrt{B}}
           \end{CD}
         \end{equation*}
         kita memperoleh bahwa matriks bentuk kanoniknya adalah $PNP^{-1}$, dengan
         \begin{equation*}
           P^{-1}
           =\rep{\identity}{B,\stdbasis_5}
           =\begin{mat}[r]
              1 &0 &0 &0 &0 \\              
              0 &1 &0 &1 &0 \\
              1 &0 &1 &0 &0 \\
              0 &0 &1 &1 &1 \\
              0 &0 &0 &0 &1
            \end{mat}
         \end{equation*}
         dan $P$ adalah invers matriks tersebut.
         \begin{equation*}
           P=\rep{\identity}{\stdbasis_5,B}
            =(P^{-1})^{-1}
           =\begin{mat}[r]
              1 &0 &0 &0 &0 \\              
             -1 &1 &1 &-1&1 \\
             -1 &0 &1 &0 &0 \\
              1 &0 &-1&1 &-1\\
              0 &0 &0 &0 &1
            \end{mat}
         \end{equation*}
        \partsitem Perhitungan untuk memeriksa hal ini bersifat rutin.
      \end{exparts}
    \end{answer}
  \recommended \item
    Setiap matriks berikut bersifat nilpoten.
    \begin{exparts*}
      \partsitem \(
        \begin{mat}[r]
          1/2  &-1/2  \\
          1/2  &-1/2
        \end{mat}        \)
      \partsitem \(
        \begin{mat}[r]
          0  &0  &0  \\
          0  &-1 &1  \\
          0  &-1 &1
        \end{mat}        \)
      \partsitem \(
        \begin{mat}[r]
         -1  &1  &-1 \\
          1  &0  &1  \\
          1  &-1 &1
        \end{mat}        \)
    \end{exparts*}
    Nyatakan masing-masing dalam bentuk kanonik.
    \begin{answer}
      \begin{exparts}
      \partsitem Perhitungan
        \begin{center}
          \begin{tabular}{c|cc}
             \multicolumn{1}{c}{\( p \)}  &\( N^p \)  &\( \nullspace{N^p} \) \\
            \hline
             \( 1 \)
               &\matrixvenlarge{\begin{mat}[r]
                    1/2  &-1/2  \\
                    1/2  &-1/2 
                  \end{mat}}
               &\( \set{\matrixvenlarge{\colvec{u \\ u}} 
                              \suchthat u\in\C}  \) \\
            \( 2 \)
               &\textit{--matriks nol--}
               &\( \C^2 \)
          \end{tabular}
        \end{center}
        menunjukkan bahwa setiap pemetaan yang direpresentasikan oleh matriks
        tersebut harus beraksi pada basis untaian dengan cara berikut.
        \begin{equation*}
          \begin{strings}{ccccccc}
            \vec{\beta}_1 &\mapsto &\vec{\beta}_2  &\mapsto &\zero  
          \end{strings}
        \end{equation*}
        karena ruang nol setelah satu kali penerapan berdimensi satu
        dan tepat satu vektor basis, $\vec{\beta}_2$, dipetakan ke nol.
        Oleh karena itu, representasi terhadap
        $\sequence{\vec{\beta}_1,\vec{\beta}_2}$ ini merupakan bentuk kanonik.
        \begin{equation*}
          \begin{mat}[r]
            0    &0   \\
            1    &0
          \end{mat}        
        \end{equation*}
      \partsitem Perhitungan di sini serupa dengan perhitungan sebelumnya.
        \begin{center}
          \begin{tabular}{c|cc}
             \multicolumn{1}{c}{\( p \)}  &\( N^p \)  &\( \nullspace{N^p} \) \\
             \hline
             \( 1 \)
               &\matrixvenlarge{\begin{mat}[r]
                    0  &0  &0  \\
                    0  &-1 &1  \\
                    0  &-1 &1  
                  \end{mat}}
               &\( \set{\matrixvenlarge{\colvec{u \\ v \\ v}} 
                              \suchthat u,v\in\C}  \) \\
           \( 2 \)
               &\textit{--matriks nol--}
               &\( \C^3 \)
          \end{tabular}
       \end{center}
       Tabel tersebut menunjukkan bahwa basis untaiannya berbentuk
        \begin{equation*}
          \begin{strings}{ccccccc}
            \vec{\beta}_1 &\mapsto &\vec{\beta}_2 &\mapsto &\zero  \\
            \vec{\beta}_3 &\mapsto &\zero
          \end{strings}
        \end{equation*}
        karena ruang nol setelah satu kali penerapan pemetaan berdimensi
        dua\Dash $\vec{\beta}_2$ dan $\vec{\beta}_3$ keduanya dipetakan ke
        nol\Dash dan satu iterasi lagi membuat vektor tambahan tersebut
        dipetakan ke nol.
      \partsitem Perhitungan
        \begin{center}
          \begin{tabular}{c|cc}
             \multicolumn{1}{c}{\( p \)}  &\( N^p \)  &\( \nullspace{N^p} \) \\
             \hline
             \( 1 \)
               &\matrixvenlarge{\begin{mat}[r]
                    -1  &1  &-1  \\
                     1  &0  &1   \\
                     1  &-1 &1 
                  \end{mat}}
               &\( \set{\matrixvenlarge{\colvec{u \\ 0 \\ -u}} 
                              \suchthat u\in\C}  \) \\
             \( 2 \)
               &\matrixvenlarge{\begin{mat}[r]
                     1  &0  &1   \\
                     0  &0  &0   \\
                    -1  &0  &-1
                  \end{mat}}
               &\( \set{\matrixvenlarge{\colvec{u \\ v \\ -u}} 
                              \suchthat u,v\in\C}  \) \\
            \( 3 \)
               &\textit{--matriks nol--}
               &\( \C^3 \)
          \end{tabular}
        \end{center}
        menunjukkan bahwa setiap pemetaan yang direpresentasikan oleh matriks
        ini harus beraksi pada basis untaian dengan cara berikut.
        \begin{equation*}
          \begin{strings}{ccccccc}
            \vec{\beta}_1 &\mapsto &\vec{\beta}_2 &\mapsto 
              &\vec{\beta}_3 &\mapsto &\zero  
          \end{strings}
        \end{equation*}
        Oleh karena itu, berikut ini adalah bentuk kanoniknya.
        \begin{equation*}
          \begin{mat}[r]
            0    &0   &0   \\
            1    &0   &0   \\
            0    &1   &0
          \end{mat}        
        \end{equation*}
      \end{exparts}  
     \end{answer}
  \item 
    Jelaskan pengaruh perkalian kiri atau kanan dengan suatu matriks yang
    berada dalam bentuk kanonik untuk matriks nilpoten.
    \begin{answer}
      Dua contoh berikut
      \begin{equation*}
        \begin{mat}[r]
          0  &0  \\
          1  &0
        \end{mat}
        \begin{mat}
          a  &b  \\
          c  &d
        \end{mat}
        =
        \begin{mat}
          0  &0  \\
          a  &b  
        \end{mat}
        \qquad
        \begin{mat}[r]
          0  &0  &0 \\
          1  &0  &0 \\
          0  &1  &0
        \end{mat}
        \begin{mat}
          a  &b  &c \\
          d  &e  &f \\
          g  &h  &i
        \end{mat}
        =
        \begin{mat}
          0  &0  &0 \\
          a  &b  &c \\
          d  &e  &f
        \end{mat}
      \end{equation*}
      menunjukkan bahwa perkalian kiri dengan suatu blok entri satu pada
      subdiagonal menggeser baris-baris matriks ke bawah.
      Blok-blok yang berbeda
      \begin{equation*}
        \begin{mat}[r]
          0  &0  &0  &0  \\
          1  &0  &0  &0  \\
          0  &0  &0  &0  \\
          0  &0  &1  &0
        \end{mat}
        \begin{mat}
          a  &b  &c  &d  \\
          e  &f  &g  &h  \\
          i  &j  &k  &l  \\
          m  &n  &o  &p
        \end{mat}
        =
        \begin{mat}
          0  &0  &0  &0  \\
          a  &b  &c  &d  \\       
          0  &0  &0  &0  \\
          i  &j  &k  &l  
        \end{mat}
      \end{equation*}
      menggeser bagian-bagian berbeda dari matriks ke bawah.

      Perkalian kanan memberikan efek yang analog pada kolom.
      Lihat \nearbyexercise{exer:IndNilLftShift}.
    \end{answer}
  \item 
    Apakah sifat nilpoten invarian terhadap keserupaan?
    Artinya, apakah matriks yang serupa dengan matriks nilpoten juga harus nilpoten?
    Jika ya, apakah indeksnya sama?
    \begin{answer}
      Ya.
      Generalisasikan kalimat terakhir dalam \nearbyexample{ex:NilMatNotCanon}.
      Mengenai indeksnya, kalimat terakhir yang sama menunjukkan bahwa indeks
      matriks baru kurang dari atau sama dengan indeks $\hat{N}$; dengan
      menukar peran kedua matriks, diperoleh pertidaksamaan dalam arah sebaliknya.

      Jawaban lain adalah menunjukkan bahwa suatu matriks
      nilpoten jika dan hanya jika setiap pemetaan terkaitnya nilpoten,
      dengan indeks yang sama.
      Karena matriks-matriks serupa merepresentasikan pemetaan yang sama,
      kesimpulan tersebut pun mengikuti.
      Inilah isi \nearbyexercise{exer:MatNilIffMapNil} di bawah.
    \end{answer}  
  \recommended \item
    Tunjukkan bahwa satu-satunya nilai eigen matriks nilpoten adalah nol.
    \begin{answer}
      Perhatikan bahwa matriks nilpoten dalam bentuk kanonik hanya mempunyai
      nilai eigen nol; misalnya, determinan matriks segitiga bawah berikut
      \begin{equation*}
         \begin{mat}
           -x  &0  &0  \\
            1  &-x &0  \\
            0  &1  &-x
         \end{mat}
      \end{equation*}
      adalah \( (-x)^3 \), yang satu-satunya akarnya adalah nol.
      Namun, matriks-matriks serupa mempunyai nilai eigen yang sama dan setiap
      matriks nilpoten serupa dengan matriks dalam bentuk kanonik.

      Cara lain untuk melihatnya adalah dengan mengamati bahwa matriks nilpoten
      membawa semua vektor ke nol setelah sejumlah iterasi, tetapi hal ini
      bertentangan dengan aksi pada ruang eigen
      $\vec{v}\mapsto \lambda\vec{v}$ kecuali jika $\lambda$ nol.
    \end{answer}
  \item 
    Adakah transformasi nilpoten berindeks tiga pada ruang
    berdimensi dua?
    \begin{answer}
      Tidak. Menurut \nearbylemma{le:RangeAndNullChains}, untuk pemetaan pada
      ruang berdimensi dua, nulitas telah mencapai nilai sebesar mungkin
      pada iterasi kedua.
    \end{answer}
  \item 
    Dalam bukti \nearbytheorem{th:NilMapHasStrBas}, mengapa kasus dasar
    buktinya bukan indeks nilpotensi nol?
    \begin{answer}
      Indeks nilpotensi suatu transformasi hanya dapat bernilai nol ketika
      vektor awal untaian harus berupa $\zero$, yakni hanya ketika $V$
      merupakan ruang trivial.
    \end{answer}
  \recommended \item
    Misalkan \( \map{t}{V}{V} \) suatu transformasi linear dan andaikan
    \( \vec{v}\in V \) sedemikian sehingga \( t^k(\vec{v})=\zero \), tetapi
    \( t^{k-1}(\vec{v})\neq\zero \).
    Tinjau untaian~$t$
    $\sequence{\vec{v},t(\vec{v}),\dots,t^{k-1}(\vec{v})}$.
    \begin{exparts}
      \partsitem Buktikan bahwa \( t \) merupakan transformasi pada rentang
        himpunan vektor dalam untaian tersebut; yakni, buktikan bahwa
        pembatasan \( t \) pada rentang itu mempunyai ruang hasil yang
        merupakan himpunan bagian rentang tersebut.
        Kita mengatakan bahwa rentang tersebut merupakan subruang
        \definend{invarian terhadap \( t \)}.\index{subruang invarian}
      \partsitem Buktikan bahwa pembatasan tersebut nilpoten.
      \partsitem Buktikan bahwa untaian~$t$ tersebut
        bebas linear, sehingga merupakan basis bagi rentangnya.
      \partsitem Representasikan pemetaan pembatasan terhadap
        basis untaian~$t$.
    \end{exparts}
    \begin{answer}
      \begin{exparts}
        \partsitem Setiap anggota $\vec{w}$ dari rentang tersebut merupakan
          kombinasi linear
          $\vec{w}=c_0\cdot \vec{v}+c_1\cdot t(\vec{v})+\dots
           +c_{k-1}\cdot t^{k-1}(\vec{v})$.
          Maka, berdasarkan linearitas pemetaan,
          $t(\vec{w})=c_0\cdot t(\vec{v})+c_1\cdot t^2(\vec{v})+\dots
           +c_{k-2}\cdot t^{k-1}(\vec{v})+c_{k-1}\cdot \zero$
          juga berada dalam rentang tersebut.
        \partsitem Operasi pada butir sebelumnya, bila diiterasikan $k$ kali,
          akan menghasilkan kombinasi linear vektor-vektor nol.
        \partsitem Hipotesis \(t^{k-1}(\vec v)\neq\zero\) memastikan bahwa
          \(\vec v\neq\zero\). Tuliskan
          \( c_0\vec{v}+\dots+c_{k-1}t^{k-1}(\vec{v})=\zero \)
          dan terapkan \( t^{k-1} \) pada kedua ruas.
          Ruas kanan menghasilkan \( \zero \), sedangkan ruas kiri menghasilkan
          \( c_0t^{k-1}(\vec{v}) \); simpulkan bahwa \( c_0=0 \).
          Lanjutkan dengan cara ini, dengan menerapkan \( t^{k-2} \) pada
          kedua ruas, dan seterusnya.
        \partsitem Tentu saja, $t$ beraksi pada rentang tersebut dengan beraksi
          pada basis ini sebagai satu untaian~$t$ berpanjang~$k$.
          \begin{equation*}
            \begin{mat}
              0 &0 &0 &0 &\ldots &0 &0 \\
              1 &0 &0 &0 &\ldots &0 &0 \\
              0 &1 &0 &0 &\ldots &0 &0 \\
              0 &0 &1 &0 &       &0 &0 \\
                &  &  &\ddots          \\
              0 &0 &0 &0 &       &1 &0 \\
            \end{mat}
          \end{equation*}
      \end{exparts}
    \end{answer}
  \item \label{exer:NilMapWillHaveStrBas}
    Selesaikan bukti \nearbytheorem{th:NilMapHasStrBas}.
    \begin{answer}
      Kita harus memeriksa bahwa
      \( B\union\set{\vec{\xi}_1,\dots,\vec{\xi}_p}
         \union\set{\vec{v}_1,\dots,\vec{v}_i} \) bebas linear,
      dengan \( B \) basis untaian \( t \) bagi \( \rangespace{t} \),
      dengan
      \( \hat C=\cat{C}{\sequence{\vec\xi_1,\dots,\vec\xi_p}} \)
      basis bagi \( \nullspace{t} \), dan dengan
      \( t(\vec{v}_1)=\vec{\beta}_1,\dots,t(\vec{v}_i)=\vec{\beta}_i \).
      Tuliskan
      \begin{equation*}
        \zero=
        \sum_{r=1}^{i}\left(c_{r,-1}\vec v_r
          +\sum_{s=0}^{h_r}c_{r,s}t^s(\vec\beta_r)\right)
          +\sum_{q=1}^{p}d_q\vec\xi_q
      \end{equation*}
      lalu terapkan \( t \).
      \begin{multline*}
        \zero=
        \sum_{r=1}^{i}\left(c_{r,-1}\vec\beta_r
          +\sum_{s=0}^{h_r-1}c_{r,s}t^{s+1}(\vec\beta_r)\right),
      \end{multline*}
      Karena \(B\) adalah basis, semua
      \(c_{r,-1},c_{r,0},\dots,c_{r,h_r-1}\) bernilai nol.
      Substitusi kembali ke persamaan pertama menyisakan kombinasi linear
      vektor-vektor dalam \(\hat C\). Karena \(\hat C\) adalah basis,
      semua \(c_{r,h_r}\) dan \(d_q\) juga nol.
    \end{answer}
  \item \label{exer:MatNilIffMapNil} 
    Tunjukkan bahwa istilah `transformasi nilpoten' dan `matriks nilpoten',
    sebagaimana diberikan dalam \nearbydefinition{def:nilpotent}, saling
    bersesuaian:~suatu pemetaan nilpoten jika dan hanya jika direpresentasikan
    oleh matriks nilpoten.
    (Apakah suatu transformasi nilpoten jika dan hanya jika ada basis yang
    terhadapnya representasi pemetaan merupakan matriks nilpoten, ataukah
    setiap representasinya merupakan matriks nilpoten?)
    \begin{answer}
      Untuk sebarang basis $B$,
      transformasi~$n$ nilpoten jika dan hanya jika
      $N=\rep{n}{B,B}$ merupakan matriks nilpoten.
      Hal ini karena hanya matriks nol yang merepresentasikan pemetaan nol,
      sehingga \( n^j \) merupakan pemetaan nol jika dan hanya jika \( N^j \)
      merupakan matriks nol.
    \end{answer}
  \item 
    Misalkan \( T \) nilpoten dengan indeks empat.
    Seberapa besar dimensi ruang hasil \( T^3 \) yang mungkin?
    \begin{answer}
      Dimensinya dapat berupa sebarang bilangan yang lebih besar dari atau
      sama dengan satu.
      Untuk memperoleh transformasi nilpoten berindeks empat yang kubiknya
      mempunyai ruang hasil berdimensi~$k$, ambil suatu ruang vektor,
      suatu basis bagi ruang tersebut, dan suatu transformasi yang beraksi
      pada basis itu dengan cara berikut.
      \begin{equation*}
         \begin{strings}{ccccccccc}
            \vec{\beta}_1 &\mapsto &\vec{\beta}_2 &\mapsto &\vec{\beta}_3
              &\mapsto &\vec{\beta}_4 &\mapsto &\zero  \\
            \vec{\beta}_5 &\mapsto &\vec{\beta}_6 &\mapsto &\vec{\beta}_7
              &\mapsto &\vec{\beta}_8 &\mapsto &\zero  \\
            &\vdotswithin{\vec{\beta}_{4k-3}}                       \\
            \vec{\beta}_{4k-3} &\mapsto &\vec{\beta}_{4k-2} 
              &\mapsto &\vec{\beta}_{4k-1}
              &\mapsto &\vec{\beta}_{4k} &\mapsto &\zero  \\
            &\vdotswithin{\vec{\beta}_{4k-3}}              \\
            \multicolumn{8}{l}{%
              \text{\textit{--mungkin ada untaian lain yang lebih pendek--}}}
         \end{strings}
      \end{equation*}
      Jadi, dimensi ruang hasil $T^3$ dapat sebesar yang diinginkan.
      Nilai terkecilnya adalah satu\Dash harus ada sedikitnya satu untaian;
      jika tidak, indeks nilpotensi pemetaan tersebut tidak akan bernilai empat.
    \end{answer}
  \item 
    Ingat bahwa matriks-matriks serupa mempunyai nilai eigen yang sama.
    Tunjukkan bahwa konversnya tidak berlaku.
    \begin{answer}
      Kedua matriks berikut hanya mempunyai nol sebagai nilai eigen:
      \begin{equation*}
        \begin{mat}
          0  &0  \\
          0  &0
        \end{mat}
        \qquad
        \begin{mat}
          0  &0  \\
          1  &0
        \end{mat}
      \end{equation*}
      tetapi keduanya tidak serupa (wakil kanoniknya berbeda, yakni
      masing-masing matriks itu sendiri).
    \end{answer}
  \item \nearbylemma{lem:RestONeToOne} menunjukkan bahwa untuk sebarang
    transformasi linear \( \map{t}{V}{V} \), pembatasan
    \( \map{t}{\genrangespace{t}}{\genrangespace{t}} \)
    bersifat satu-satu.
    Tunjukkan bahwa pembatasan itu juga pada, sehingga merupakan automorfisme.
    Haruskah pembatasan tersebut berupa pemetaan identitas?
    \begin{answer}
      Pemetaan itu pada menurut \nearbylemma{le:RangeAndNullChains}.
      Pemetaan tersebut tidak harus berupa identitas; tinjau pemetaan
      $\map{t}{\Re^2}{\Re^2}$ berikut.
      \begin{equation*}
        \colvec{x \\ y}\mapsunder{t}\colvec{y \\ x}
      \end{equation*}
      Untuk pemetaan ini $\genrangespace{t}=\Re^2$, dan $t$ bukan identitas.
    \end{answer}
  \item 
    Buktikan bahwa suatu matriks nilpoten serupa dengan matriks yang seluruh
    entrinya nol, kecuali blok-blok entri satu pada superdiagonal.
    \begin{answer}
      Cukup lakukan pengurutan ulang sederhana pada basis untaian.
      Sebagai contoh, suatu pemetaan yang terkait dengan basis untaian berikut
      \begin{equation*}
         \begin{strings}{ccccccc}
            \vec{\beta}_1 &\mapsto &\vec{\beta}_2 &\mapsto &\zero  
         \end{strings}
      \end{equation*}
      direpresentasikan terhadap
      $B=\sequence{\vec{\beta}_1,\vec{\beta}_2}$ oleh matriks berikut,
      \begin{equation*}
        \begin{mat}
          0  &0  \\
          1  &0          
        \end{mat}
      \end{equation*}
      tetapi direpresentasikan terhadap
      $B=\sequence{\vec{\beta}_2,\vec{\beta}_1}$ dengan cara berikut.
      \begin{equation*}
        \begin{mat}
          0  &1  \\
          0  &0          
        \end{mat}
      \end{equation*}
    \end{answer}
  \recommended \item
    Buktikan bahwa jika suatu transformasi mempunyai ruang hasil yang sama
    dengan ruang nolnya, maka dimensi domainnya genap.
    \begin{answer}
      Misalkan $\map{t}{V}{V}$ adalah transformasi tersebut.
      Jika \( \rank (t)=\nullity (t) \), maka persamaan
      \( \rank(t)+\nullity(t)=\dim (V) \) menunjukkan bahwa
      \( \dim (V) \) genap.
    \end{answer}
  \item 
    Buktikan bahwa jika dua matriks nilpoten saling komutatif, maka hasil kali
    dan jumlah keduanya juga nilpoten.
    \begin{answer}
      Agar nilpoten, kedua matriks harus persegi.
      Agar saling komutatif, keduanya harus berukuran sama.
      Jadi, hasil kali dan jumlahnya terdefinisi.

      Namai kedua matriks itu \( A \) dan \( B \).
      Untuk melihat bahwa \( AB \) nilpoten, hitung
      $
         (AB)^2=ABAB=AABB=A^2B^2$,
      serta
         $(AB)^3=A^3B^3$, dan seterusnya.
      Karena \( A \) nilpoten, hasil kali tersebut pada akhirnya bernilai nol.

      Untuk jumlah, argumennya serupa; gunakan Teorema Binomial.
     \end{answer}
  % 2014-Dec-19 These two questions and answers are fine; I commented them to make the pages fit better.
  % \item 
  %   Consider the transformation of \( \matspace_{\nbyn{n}} \) given
  %   by \( t_S(T)=ST-TS \) where \( S \) is an \( \nbyn{n} \) matrix.
  %   Prove that if \( S \) is nilpotent then so is \( t_S \).
  %   \begin{answer}
  %     Some experimentation gives the idea for the proof.
  %     Expansion of the second power
  %     \begin{equation*}
  %       t^2_S(T)=S(ST-TS)-(ST-TS)S=S^2-2STS+TS^2
  %     \end{equation*}
  %     the third power
  %     \begin{align*}
  %       t^3_S(T)
  %         &=S(S^2-2STS+TS^2)-(S^2-2STS+TS^2)S  \\
  %         &=S^3T-3S^2TS+3STS^2-TS^3
  %     \end{align*}
  %     and the fourth power
  %     \begin{align*}
  %       t^4_S(T)
  %         &=S(S^3T-3S^2TS+3STS^2-TS^3)-(S^3T-3S^2TS+3STS^2-TS^3)S  \\
  %         &=S^4T-4S^3TS+6S^2TS^2-4STS^3+TS^4
  %     \end{align*}
  %     suggest that the expansions follow the Binomial Theorem.
  %     Verifying this by induction on the power of $t_S$ is routine.
  %     This answers the question because, where the index of nilpotency of 
  %     $S$ is $k$, in the expansion of $t^{2k}_S$  
  %     \begin{equation*}
  %       t^{2k}_S(T)=\sum_{0\leq i\leq 2k}(-1)^i\binom{2k}{i} S^iTS^{2k-i}
  %     \end{equation*}
  %     for any $i$ at least one of the $S^i$ and $S^{2k-i}$ 
  %     has a power higher than $k$, and so the term gives the zero matrix. 
  %   \end{answer}
  % \item 
  %   Show that if \( N \) is nilpotent then \( I-N \) is
  %   invertible.
  %   Is that `only if' also?
  %   \begin{answer}
  %     Use the geometric series:
  %     $
  %        I-N^{k+1}=(I-N)(N^k+N^{k-1}+\cdots+I)
  %     $.
  %     If \( N^{k+1} \) is the zero matrix then we have a right inverse for
  %     \( I-N \).
  %     It is also a left inverse.

  %     This statement is not `only if' since
  %     \begin{equation*}
  %        \begin{mat}[r]
  %           1  &0  \\
  %           0  &1
  %        \end{mat}
  %        -\begin{mat}[r]
  %           -1  &0  \\
  %           0  &-1
  %        \end{mat}
  %     \end{equation*}
  %     is invertible.  
  %    \end{answer}
%  \recommended \item 
%    Show that when a nilpotent matrix \( T \) is in canonical form 
%    then the number of
%    blocks of \( \nbyn{(k+1)} \) matrices that are all zeros
%    except for subdiagonal ones is
%    \( 2\,(\nullity (T^k))
%      -\nullity (T^{k+1})-\nullity (T^{k-1}) \).
\index{nilpoten|)}
\end{exercises}
